\section{Absorbers: forward patterns and early cut-off}\label{sec:absorbedores}

Fix $1\le h\le K$ and work under $\QQ_h$. We track the queens that could be removed to make a repair at a given crossing. The degree
estimates from the preceding section allow us to estimate how their number changes when the next piece is placed.

\subsection{Repairs at a crossing and line marks}

The \emph{root} is the square $(r,c)$ that we track as a possible empty crossing. In Figure~\ref{fig:reparacion-tablero} it would be
$(2,1)$. The candidate queen occupies $(0,0)$ and its two \emph{destinations} are $(2,0)$ and $(0,1)$. We first check that the destinations are
compatible with each other. We call this condition the $G$ filter. We then check that the diagonals of both destinations are free of other
queens.

Fix a root $(r,c)$. Let $\mathcal I_{rc}$ be the set of squares $e=(i,j)$ with $i\ne r$ and $j\ne c$ whose \emph{destinations} $(r,j)$ and
$(i,c)$ do not attack each other (filter $G=1$). This excludes at most $Cn$ squares. The \emph{candidates} for the root are the
squares in $\mathcal I_{rc}$. Candidates that contain a queen are called \emph{occupied}. The four lines of $e$ and the four diagonals
and antidiagonals of its destinations are eight distinct vertices. This is where we use the $G$ filter. We denote the latter four lines by
$d_1,\dots,d_4$ and their list by $\mathcal D(r,c;i,j)$.

For a vertex $v$, let $\chi_v$ indicate that $v$ is active and let $V_v$ indicate that line $v$ is physically free of queens. We always
have $\chi_v\le V_v$. The \emph{mark} of a line $v$ is its state: the line carries the active mark if $\chi_v=1$ and the mark $V$ if $V_v=1$. Let
$O_e$ indicate that a real queen has been placed at $e$, and let $\chi_e=\prod_{v\in e}\chi_v$. For $J\subseteq\{1,2,3,4\}$ set
\begin{equation}\label{eq:patrones}
 M_J=\prod_{i\in J}\chi_{d_i}\prod_{i\notin J}V_{d_i},\qquad
 Z^\oc_{e,J}=O_e\,M_J,\qquad Z^\di_{e,J}=\chi_e\,M_J .
\end{equation} There are $32$ patterns per square, and apart from the remainder terms in Lemma~\ref{lem:generador}, the process
generates no others. The extended physical counter is \[
 A^{rc}=\sum_{e\in\mathcal I_{rc}}Z^\oc_{e,\varnothing}=\sum_{e\in\mathcal I_{rc}}O_e\prod_{i=1}^4V_{d_i}.
\] If row $r$ and column $c$ are physically empty, $A^{rc}$ equals the number of legal one-move repairs at crossing $(r,c)$. Indeed,
removing the queen at $e$ and placing queens at its two destinations is legal exactly when the four destination diagonals are empty. The rescaled
weighted availability is $\overline Y^{rc}=\sum_{e\in\mathcal I_{rc}}v_eZ^\di_{e,\varnothing}$. With the original weights it is
$x_*^2\overline Y^{rc}$.

\subsection{Reference generator}

In discrete time, with $z=(N-t)/N$, define \[
 \alpha=\frac1{Nz},\qquad \beta_i=\frac{p_{d_i}}{Nz},\qquad k_e=\frac{\pi_e}{Nz^4},
\] and the linear operator \begin{equation}\label{eq:generador}
\begin{aligned}
 \cL Z^\oc_{e,J}&=-|J|\,\alpha\,Z^\oc_{e,J}-\sum_{i\notin J}\beta_i\,Z^\oc_{e,J\cup\{i\}}+k_e\,Z^\di_{e,J},\\
 \cL Z^\di_{e,J}&=-(4+|J|)\,\alpha\,Z^\di_{e,J}-\sum_{i\notin J}\beta_i\,Z^\di_{e,J\cup\{i\}} .
\end{aligned}
\end{equation} An active mark dies at rate $\alpha$ when hit by a real or virtual piece. A mark $V$ dies only when hit by a real piece,
which requires the corresponding active mark. This turns $V_{d_i}$ into a contribution with $\chi_{d_i}$ at rate $\beta_i$. A virtual
death affects $\chi$ but leaves $V$ intact. The birth probability for $O_e$ is exactly $[q_t/(N-t)]\,v_e/W$ times the indicator that
its four lines are active. In the region~\eqref{eq:regularidad} this probability is $k_e\chi_e$ up to error terms. The birth does not
affect the four destination lines, which are distinct from its own lines.

\begin{lema}[Comparison with the true generator]\label{lem:generador} In the region~\eqref{eq:regularidad}, the true drifts differ
from~\eqref{eq:generador} by remainder terms satisfying \[
 |R^\oc_{e,J}|\le C_\zeta\bigl[(\eps/n+n^{-2})\,O_e+\eps n^{-2}\,Z^\di_{e,J}\bigr],\qquad
 |R^\di_{e,J}|\le C_\zeta(\eps/n+n^{-2}).
\] \end{lema}

\begin{proof} Adjusting $\eps$ absorbs the error in the colors from~\eqref{eq:regularidad}. The rates~\eqref{eq:tasas} control each
single deletion, and the probability of deleting two distinct marks in one step is $O_\zeta(n^{-2})$. Selecting $e$ deletes its four
active marks. This is included in the error for $Z^\di$, but its birth in $Z^\oc$ is retained exactly. \end{proof}

We do not factor any joint probability under $\mu_h$. The generator acts on indicators of the reference process.
